\documentclass[11pt,a4paper,sans]{moderncv}

% Style options
\moderncvstyle{classic}        % casual / classic / banking / fancy
\moderncvcolor{black}          % black, blue, green, red, purple, orange, grey

% Encoding & fonts (IMPORTANT: use XeLaTeX)
\usepackage[T1]{fontenc}
\usepackage{lmodern}

% Layout tweaks
\usepackage[scale=0.75]{geometry}

% Personal data
\name{}{Nikolaus Knop}
%\homepage{www.schleifenkauz.com}
%\extrainfo{Born: 2003}

\begin{document}
\makecvtitle

%-------------------------
\section{Profile}
Nikolaus Knop is a Munich-based pianist, composer, and programmer
exploring novel modes of interaction between 
live performers and computer systems --
and between acoustic and electronic sounds.
His compositional work is closely intertwined 
with the conception and implementation of original digital tools for composition and performance,
leveraging a rigorous background in computer science. 
His work as a pianist -- with a focus on the 20th-century canon,
including the music of Bartók, Schönberg, Messiaen, and Ligeti --
informs his compositional perspective.

%-------------------------
\section{Education}
\cventry{2023--}{B.Mus Piano}{University of Music and Theatre}{Munich}{}{
    \begin{itemize}
        \item     Spatialized audio performance of Bartok's \textit{The Night's Music}
        \item     Composition seminars with Hendrik Ajax and Vladimir Tarnopolski
    \end{itemize}
}
\cventry{2021--2025}{B.A. Philosophy}{LMU}{Munich}{}{
    \begin{itemize}
        \item     Led the tutorial for the lecture on philosophy of language
        \item     Bachelor thesis on musical metaphors in late Wittgenstein
    \end{itemize}
}

\cventry{2021--2023}{B.A. Computer Science}{TU}{Munich}{}{
    Studies interrupted to focus on artistic projects
}
\cventry{2018--}{Studies in Music Theory and Composition}{}{with Thomas Taxus Beck}{}{}

%-------------------------
\section{Selected Works}
\cvitem{2026}{\textit{Fantasie und Variationen} for two violins and live electronics}
\cvitem{2025}{\textit{Rondo} for clarinet, string trio, conductor, and live electronics}
\cvitem{2025}{\textit{Die Radaddistenmaschine} for fixed media}
\cvitem{2024}{\textit{FM96} for violin, piano, valve radio, and live electronics}
\cvitem{2023}{\textit{Frost} for fixed media}
\cvitem{2021}{\textit{Pareidolien} for string quartet}
%\cvitem{2020}{\textit{Obertonkanon}, fixed media}
\cvitem{2019}{\textit{Reflektionen} for solo piano}

%-------------------------
\section{Refereed Conference Presentations}
\cvitem{May 2026}{\textit{Ponticello: An Interactive Conducting System for Mixed Music Performance}, International Computer Music Conference, Hamburg}

%-------------------------
\section{Prizes \& Scholarships}
\cvitem{2021\&22}{Prizes at \textit{Jugend komponiert}}
\cvitem{2020}{First Prize at the German nation computer science competition}
\cvitem{2020--}{Scholarship, \textit{Studienstiftung des deutschen Volkes}}

\end{document}